\documentclass[10pt,a4paper]{article}

% ----------------- Paquetes -------------------%
\usepackage[T1]{fontenc}
\usepackage[utf8]{inputenc}
\usepackage[a4paper,margin=2.5cm]{geometry}
\usepackage{mathptmx}             
\usepackage{changepage} 
\usepackage{setspace}
\usepackage[spanish]{babel}
\usepackage[labelfont=bf,labelsep=space]{caption}
\usepackage{amssymb}
\usepackage{amsmath}
\usepackage{ebgaramond}
\usepackage{placeins}
\usepackage{etoolbox}
\usepackage{setspace}
\usepackage{parskip} 
\usepackage{tcolorbox}
\usepackage{needspace}
\usepackage{xparse}
\usepackage{ocgx2}
\usepackage{hyperref}
\usepackage{hypcap}


%%%%%%%%%%%%%%%%%%%%%%%%%%%%%%%%%%%%%%%%%%%%%%%%%%%%%%%%%%%%%%%%
% --- Fija primero tus valores globales "buenos" ---
\setstretch{1.2}                 % Interlineado global
\setlength{\parskip}{0.7\baselineskip}
\setlength{\parindent}{0pt}
\newlength{\GlobalParskip}
\setlength{\GlobalParskip}{\parskip}

\newcounter{eqnum}[subsection]
\renewcommand{\theeqnum}{\arabic{eqnum}}
\newcounter{alphaitem}
\newcounter{numitem}

%% ------------------- Recuadros----------------------%%
\newcommand{\puntosclave}[1]{%
  \begin{tcolorbox}[
    colframe=black,
    title={Puntos clave},
    before upper={\setlength{\parindent}{0.5cm}\setlength{\parskip}{0.5\baselineskip}}
  ]
  #1
  \end{tcolorbox}
}

\newcommand{\modificaciones}[1]{
  \begin{tcolorbox}[
    colback=white,
    colframe=red,
    title={Modificaciones a realizar},
    before upper={\setlength{\parindent}{0.5cm}\setlength{\parskip}{0.5\baselineskip}}
  ]
  #1
  \end{tcolorbox}
}

%% ------------------- Preguntas TEST----------------------%%
% Opciones: radio (single-choice)
\NewDocumentCommand{\OptRadio}{m m m}{%
  \noindent\CheckBox[radio,name=#1,value=#2,width=1.2ex,height=1.2ex]{}%
  \;\textbf{#2.}~#3\par
}

% Opciones: checkbox (multi-choice)
\NewDocumentCommand{\OptCheck}{m m}{%
  \noindent\CheckBox[name=#1,width=1.2ex,height=1.2ex,bordercolor={0 0 0}]{}%
  \;#2\par
}

% Pregunta single-choice (radio)
% Args: 1=id, 2=enunciado, 3=opciones, 4=solución+justificación, 5=imagen (o {}), 6=origen
\NewDocumentCommand{\PreguntaTestSingle}{m m m m m m}{%
  \par\medskip
  \noindent\textbf{#1.}~#2\par\smallskip

    % Imagen (si no es {})
    \IfBlankTF{#5}{}{%
      \begin{center}
        \includegraphics[width=0.75\linewidth]{#5}
      \end{center}
    }

    \begin{Form}
      #3
    \end{Form}

    \par\smallskip
    \switchocg{sol-#1}{\fbox{Mostrar / ocultar solución}}%
    \hfill{\footnotesize #6}\par\smallskip

    \begin{ocg}{Solución}{sol-#1}{0}
      \noindent #4
    \end{ocg}
  \par\medskip
}

% Pregunta multi-choice (checkboxes)
\NewDocumentCommand{\PreguntaTestMulti}{m m m m m m}{%
  \par\medskip
  \noindent\textbf{#1.}~#2\par\smallskip

    % Imagen (si no es {})
    \IfBlankTF{#5}{}{%
      \begin{center}
        \includegraphics[width=0.75\linewidth]{#5}
      \end{center}
    }
    \begin{Form}
      #3
    \end{Form}

    \par\smallskip
    \switchocg{sol-#1}{\fbox{Mostrar / ocultar solución}}%
    \hfill{\footnotesize #6}\par\smallskip

    \begin{ocg}{Solución}{sol-#1}{0}
      \noindent #4
    \end{ocg}
  \par\medskip
}

%% ------------------- Listas----------------------%%
    % --- Lista alfabética ---
    \newenvironment{la}
      {\begin{list}{\alph{alphaitem})}{%
          \usecounter{alphaitem}%
          \setlength{\leftmargin}{1cm}%
          \setlength{\parskip}{3pt}% 
          \setlength{\itemsep}{0pt}% ← espacio entre ítems
          \setlength{\parsep}{0pt}%  ← párrafos dentro de un ítem
      }}
      {\end{list}}
      
    % --- Lista numérica ---
    \newenvironment{lnum}
      {\begin{list}{\arabic{numitem}.}{%
          \usecounter{numitem}%
          \setlength{\leftmargin}{1cm}%
          \setlength{\parskip}{3pt}% 
          \setlength{\itemsep}{0pt}% ← espacio entre ítems
          \setlength{\parsep}{0pt}%  ← párrafos dentro de un ítem
      }}
      {\end{list}}
      
% ------------------- Imágenes----------------------%
    \graphicspath{{Img/}}% Rutas por defecto 
    % --- Imagen adaptativa ---
    % Registros de dimensión definidos una sola vez
    \newdimen\imgcap
    \newdimen\imgmin
    \newdimen\imgmax
    \newdimen\imgremain
    \newdimen\imgh
    \newdimen\imgneed
    
    \newcommand{\imagenfit}[3]{%
      \addvspace{10pt}% separación antes
      \par\begingroup
        % Parámetros de composición (asignaciones, no \newdimen)
        \imgcap=1\baselineskip            % alto estimado de título+fuente
        \imgmin=0.15\textheight
        \imgmax=0.15\textheight
        \imgremain=\pagegoal
        \ifdim\pagegoal=\maxdimen \imgremain=0pt\fi
        \advance\imgremain by -\pagetotal
        \advance\imgremain by -\imgcap
        % Decisión de altura destino
        \ifdim\imgremain<\imgmin
          \newpage
          \imgh=0.20\textheight
        \else
          \imgh=\imgremain
          \ifdim\imgh>\imgmax \imgh=\imgmax \fi
        \fi
        % Reservar espacio para evitar cortes del bloque completo
        \imgneed=\imgh \advance\imgneed by \imgcap
        \Needspace{\imgneed}
        % Cuerpo
        \noindent\begin{minipage}{\textwidth}
          \centering
          \captionof{figure}{\textemdash\ #2}\par\vspace{0.3em}% título arriba
          \includegraphics[width=\linewidth,height=\imgh,keepaspectratio]{\detokenize{#1}}\par
          \vspace{0.3em}\small\textit{Fuente: }#3% fuente abajo
        \end{minipage}%
      \endgroup
      \par\addvspace{10pt}%
    }

%------------ Bloque de RECUADRO ESPECIAL -------------%
    \newenvironment{specialblock}[1]{%
      \begin{quote} % crea sangría a izquierda y derecha
      \small % texto más pequeño
      \noindent\textbf{#1}\\[-0.3em] % título en negrita
      \rule{\linewidth}{0.4pt}\\[0.6em]}{
      \end{quote}}
      
%-------------------------- Portada -----------------------%
\newcommand{\oldtitle}[1]{\def\OldTitle{#1}}
\newcommand{\changerationale}[1]{\def\ChangeWhy{#1}}

\newcommand{\changebox}{%
\begin{tcolorbox}[title={Historial de cambios del tema}]
\textbf{Título anterior:} \OldTitle\par
\textbf{Motivo del cambio:} \ChangeWhy
\end{tcolorbox}
}

%-------------------------- Author Font -----------------------%
    \newcommand{\authorfont}[1]{{\fontfamily{EBGaramond-TLF}\selectfont #1}}

%-------------------------- Anexos -----------------------%
    \makeatletter
    \newcounter{anexo}
    \renewcommand{\theanexo}{\Roman{anexo}}
    \newcommand{\anexo}[1]{%
      \clearpage
      \phantomsection
      \refstepcounter{anexo}%
      \noindent\textbf{Anexo \theanexo.- }#1\par\medskip
      \addcontentsline{stc}{section}{Anexo \theanexo.- #1}%
    }
    \makeatother

    %-------------------------- Lista de figuras (personal) -----------------------%
\makeatletter
\newcommand{\MyListOfFigures}{%
  \clearpage
  \phantomsection
  {\centering\bfseries\Large Índice de figuras\par}%
  \medskip
  % Entrada en nuestro índice de contenido (stc)
  \addcontentsline{stc}{section}{Índice de figuras}%
  % Recuperación del listado (sin cabecera propia)
  \begingroup
    \parindent=0pt\parskip=2pt
    \@starttoc{lof}%
  \endgroup
}
\makeatother

%-------------------------- Bibliografía -----------------------%
\makeatletter
% Contador para las referencias
\newcounter{myref}
% Cabecera de la bibliografía, entra en el índice 'stc'
\newcommand{\MyBibliography}{%
  \clearpage
  \phantomsection
  {\centering\bfseries\Large Bibliografía y referencias\par}%
  \medskip
  \addcontentsline{stc}{section}{Bibliografía y referencias}%
  \setcounter{myref}{0}%
}
\newcommand{\MyBibItem}[4]{%
  \refstepcounter{myref}%
  \noindent[\themyref]\ #1.\ \emph{#2}. #3. #4\par
}
\makeatother

%--------- Establecer cuatro niveles de índice --------%
    \setcounter{tocdepth}{4}   
    \setcounter{secnumdepth}{4}% numera hasta \paragraph

%%%%%%%%%%%%%%%%%%%%%%%%%%%%%%%%

\makeatletter

    %--------- Interlineado Global --------%
    \edef\GlobalBaseLineStretch{\baselinestretch}
    
    %--------- Espaciado de seccinones --------%
    \renewcommand\section{\@startsection{section}
      {1}{\z@}
      {2.5ex \@plus 1ex \@minus .2ex} % espacio antes
      {1.5ex \@plus .2ex}             % espacio después
      {\normalfont\Large\bfseries}    % formato del título
    }
    \renewcommand\subsection{\@startsection{subsection}{2}{\z@}
      {3ex \@plus 0.8ex \@minus .2ex}
      {1.5ex \@plus .2ex}
      {\normalfont\large\bfseries}
    }
    \renewcommand\subsubsection{\@startsection{subsubsection}{3}{\z@}
      {3ex \@plus 0.5ex \@minus .2ex}
      {1ex \@plus .2ex}
      {\normalfont\normalsize\bfseries}
    }
    \renewcommand\paragraph{\@startsection{paragraph}{4}{\z@}%
    {-2ex \@plus -0.5ex \@minus -.2ex}% espacio antes
    {0.8ex \@plus .2ex}% espacio después
    {\normalfont\normalsize\itshape}}% ← cursiva en lugar de negrita

    %--------- Ecuaciones --------%   
    % Variables globales para almacenar títulos y notas
    \gdef\leq@current@section{}%
    \gdef\leq@current@subsection{}%
    \gdef\leq@last@printed@section{}%
    \gdef\leq@last@printed@subsection{}%
    
    % Comando para añadir nota (invisible en el cuerpo)
    \newcommand{\eqnote}[1]{%
      \addcontentsline{leq}{leqnote}{#1}%
    }
    
    % Comando para ecuaciones
    \newcommand{\eqblock}[2]{%
      % Verificar si necesitamos imprimir el título de sección
      \ifx\leq@current@section\leq@last@printed@section\else
        \addcontentsline{leq}{leqsection}{\leq@current@section}%
        \global\let\leq@last@printed@section\leq@current@section
        \gdef\leq@last@printed@subsection{}% Reset subsección al cambiar sección
      \fi
      % Verificar si necesitamos imprimir el título de subsección
      \ifx\leq@current@subsection\leq@last@printed@subsection\else
        \ifx\leq@current@subsection\@empty\else
          \addcontentsline{leq}{leqsubsection}{\leq@current@subsection}%
          \global\let\leq@last@printed@subsection\leq@current@subsection
        \fi
      \fi
      % Ecuación
      \refstepcounter{eqnum}%
      \begin{equation*}
        #1\tag{E.\theeqnum}
      \end{equation*}%
      \addcontentsline{leq}{leqt}{[E.\theeqnum] -- #2}%
      \addcontentsline{leq}{leqm}{#1}%
    }
    
    %%% ----- Índice de ecuaciones ----------%%%
    % Tamaño para []
    \newcommand{\leqIndexFont}{\normalsize}
    \newcommand{\leqTitleFont}{\tiny}
    \newcommand{\leqNoteFont}{\tiny}
    % Cabecera de sección 
    \newcommand*\l@leqsection[2]{%
      \addvspace{25pt}% Mayor separación antes de nueva sección
      \noindent\leqIndexFont
      \makebox[\linewidth]{\rule{.38\linewidth}{0.4pt}\ \ \textbf{  \MakeUppercase{#1}}\ \ \rule{.38\linewidth}{0.4pt}}%
      \par\vspace{-10pt}
    }
    % Cabecera de subsección 
    \newcommand*\l@leqsubsection[2]{%
      \par\addvspace{15pt}%
      \noindent\leqIndexFont #1
      \par\vspace{-7pt}
    }
    % Título de ecuación 
    \newcommand*\l@leqt[2]{%
    \par\addvspace{10pt}%
      \par\noindent\leqTitleFont #1\par
    }
    % Miniatura de ecuación
    \newcommand*\l@leqm[2]{%
      \vspace{0.3\baselineskip}%
      \par\scriptsize\noindent\hspace*{1.5em}\(#1\)\par%
    }
    % Nota de ecuación 
    \newcommand*\l@leqnote[2]{%
      \vspace{0.2\baselineskip}%
      \par\noindent\leqNoteFont\hspace*{1.5em}\textbf{Nota.--} #1\par\vspace{0.5\baselineskip}%
    }
    % Generación del índice
    \newcommand{\mylistofequations}{%
      \clearpage
      \phantomsection
      % Título visual de la lista (ajústalo a tu gusto):
      {\centering\bfseries\Large Índice de ecuaciones\par}
      % Entrada en nuestro índice 'stc':
      \addcontentsline{stc}{section}{Índice de ecuaciones}%
      % Llamada a tu lista real (usa el nombre exacto de tu macro):
      \@starttoc{leq}%
    }
    
    %--------- Captura de títulos de sección --------%
    \let\leq@original@section\section
    \let\leq@original@subsection\subsection
    % Flag para saber si estamos en una sección asterisco
    \newif\ifLeqStarSection
    \LeqStarSectionfalse
    
    % Redefinir \section CON CONTROL DE ASTERISCO
    \renewcommand{\section}{%
      \@ifstar{\leq@section@star}{\@dblarg\leq@section@nostar}%
    }
    \def\leq@section@star#1{%
      \global\LeqStarSectiontrue
      \gdef\leq@current@section{#1}%
      \gdef\leq@current@subsection{}%
      \leq@original@section*{#1}%
      \global\LeqStarSectionfalse
    }
    \def\leq@section@nostar[#1]#2{%
      \global\LeqStarSectionfalse
      \gdef\leq@current@section{#2}%
      \gdef\leq@current@subsection{}%
      \leq@original@section[#1]{#2}%
    }
    % Redefinir \subsection
    \renewcommand{\subsection}{%
      \@ifstar{\leq@subsection@star}{\@dblarg\leq@subsection@nostar}%
    }
    \def\leq@subsection@star#1{%
      \global\LeqStarSectiontrue
      \gdef\leq@current@subsection{#1}%
      \leq@original@subsection*{#1}%
      \global\LeqStarSectionfalse
    }
    \def\leq@subsection@nostar[#1]#2{%
      \global\LeqStarSectionfalse
      \gdef\leq@current@subsection{#2}%
      \leq@original@subsection[#1]{#2}%
    }

    % ----- Restauración de valores Globales de estilo ----------%
    % Flags
    \newif\ifInFootnote
    \InFootnotefalse
    \pretocmd{\@footnotetext}{\InFootnotetrue}{}{}
    \apptocmd{\@footnotetext}{\InFootnotefalse}{}{}
    % Profundidad de listas (soporta anidadas)
    \newcount\MyListDepth \MyListDepth=0
    \AtBeginEnvironment{la}{\advance\MyListDepth by 1}
    \AfterEndEnvironment{la}{\advance\MyListDepth by -1}
    \AtBeginEnvironment{lnum}{\advance\MyListDepth by 1}
    \AfterEndEnvironment{lnum}{\advance\MyListDepth by -1}
    \AtBeginEnvironment{itemize}{\advance\MyListDepth by 1}
    \AfterEndEnvironment{itemize}{\advance\MyListDepth by -1}
    \AtBeginEnvironment{enumerate}{\advance\MyListDepth by 1}
    \AfterEndEnvironment{enumerate}{\advance\MyListDepth by -1}
    % Restauración diferida: hazlo al INICIO del próximo párrafo real
    \newif\if@pendingRestore
    \@pendingRestorefalse
    \newcommand{\RequestRestore}{\global\@pendingRestoretrue}
    \everypar{%
      \if@pendingRestore
        \global\@pendingRestorefalse
        \ifInFootnote\else
          \ifnum\MyListDepth=0
            \setlength{\parskip}{\GlobalParskip}%
            \setstretch{\GlobalBaseLineStretch}%
          \fi
        \fi
      \fi
    }
    % Al final de bloques:
    \AfterEndEnvironment{equation}{\RequestRestore}
    \AfterEndEnvironment{equation*}{\RequestRestore}
    \AfterEndEnvironment{la}{\RequestRestore}
    \AfterEndEnvironment{lnum}{\RequestRestore}
    % Al final de macros:
    \apptocmd{\eqblock}{\RequestRestore}{}{}
    \apptocmd{\imagenfit}{\RequestRestore}{}{}
    
\makeatother

% ===================== ToC propio con 4 niveles (único bloque) =====================
\makeatletter

% --- Escritores: añadimos una línea al .stc en cada cabecera numerada ---
% Guardamos las versiones actuales para llamar al comportamiento normal
\let\MyTOC@orig@section\section
\let\MyTOC@orig@subsection\subsection
\let\MyTOC@orig@subsubsection\subsubsection
\let\MyTOC@orig@paragraph\paragraph

% SECTION
\renewcommand{\section}{%
  \@ifstar{\MyTOC@section@star}{\@dblarg\MyTOC@section@nostar}%
}
\def\MyTOC@section@star#1{%
  \MyTOC@orig@section*{#1}% sin entrada al .stc
}
\def\MyTOC@section@nostar[#1]#2{%
  \MyTOC@orig@section[{#1}]{#2}% comportamiento normal (escribe al .toc)
  % Escribimos también nuestra línea al .stc (texto con \numberline)
  \addcontentsline{stc}{section}{\protect\numberline{\thesection}#2}%
}

% SUBSECTION
\renewcommand{\subsection}{%
  \@ifstar{\MyTOC@subsection@star}{\@dblarg\MyTOC@subsection@nostar}%
}
\def\MyTOC@subsection@star#1{%
  \MyTOC@orig@subsection*{#1}%
}
\def\MyTOC@subsection@nostar[#1]#2{%
  \MyTOC@orig@subsection[{#1}]{#2}%
  \addcontentsline{stc}{subsection}{\protect\numberline{\thesubsection}#2}%
}

% SUBSUBSECTION
\renewcommand{\subsubsection}{%
  \@ifstar{\MyTOC@subsubsection@star}{\@dblarg\MyTOC@subsubsection@nostar}%
}
\def\MyTOC@subsubsection@star#1{%
  \MyTOC@orig@subsubsection*{#1}%
}
\def\MyTOC@subsubsection@nostar[#1]#2{%
  \MyTOC@orig@subsubsection[{#1}]{#2}%
  \addcontentsline{stc}{subsubsection}{\protect\numberline{\thesubsubsection}#2}%
}

% PARAGRAPH
\renewcommand{\paragraph}{%
  \@ifstar{\MyTOC@paragraph@star}{\@dblarg\MyTOC@paragraph@nostar}%
}
\def\MyTOC@paragraph@star#1{%
  \MyTOC@orig@paragraph*{#1}%
}
\def\MyTOC@paragraph@nostar[#1]#2{%
  \MyTOC@orig@paragraph[{#1}]{#2}%
  % Si no numeras los \paragraph, cambia \theparagraph por texto plano sin \numberline
  \addcontentsline{stc}{paragraph}{\protect\numberline{\theparagraph}#2}%
}

% --- Impresor del índice propio (lee el .stc) ---
\newcommand{\MyTOC}{%
  \par\bigskip
  {\centering\bfseries\Large Índice de contenido\par}%
  \medskip
  \begingroup
    \parindent=0pt\parskip=2pt
    \@starttoc{stc}%
  \endgroup
}

\makeatother
% ===================== fin ToC propio =====================


%%%%%%%%%%%%%%


% --- Datos del tema ---
\title{La Política Comercial (II)\\
\large La política comercial estratégica. La política de promoción exterior: justificación, instrumentos y objetivos.}
\author{Marcelo Ortega}
\date{Marzo 2026}
\newcommand{\TemaCode}{3B8}

\oldtitle{
    B.8. La política comercial (II). La política comercial estratégica. La política de promoción exterior: justificación, instrumentos y objetivos.
}
\changerationale{
    Sin cambios.
}

\usepackage{soul}\sethlcolor{yellow}% resaltado amarillo: \hl{texto} (Panel TCEE, ⌃H)
\begin{document}


\maketitle
\changebox

\newpage

% --- Página de resumen y modificaciones ---
\puntosclave{ %% Escribir puntos clave Apuntes CECO
     Quedan separados claramente los temas que son de Teorías del Comercio (5 y 6) con los que son de política comercial (7 y 8). En este nuevo Tema 8 ya no hay una mención al debate teórico librecambio-proteccionismo, sino que se pide analizar directamente las políticas comerciales estratégicas y de promoción exterior que efectúan los países.\footnote{%
         La convocatoria del Martes 29 de abril de 2008 para la oposición al Cuerpo Superior de Técnicos Comerciales y Economistas del Estado modificó el título del Tema 8 de la parte B del tercer ejercicio. El nuevo terna se titula La política comercial (11): La política comercial estratégica. La política de promoción exterior: justificación, instrumentos y objetivos. El título antiguo era Lapolítica comercial (11): Librecambio y proteccionismo. La política comercial estratégica.
     } Es decir, toda vez que el Tema 7 incluye los aspectos más tradicionales (y más en desuso) de la política comercial, el Tema 8 incluye los desarrollos más modernos en PC y más utilizados en los países desarrollados, especialmente en el caso de la promoción exterior. 
     
     La optimalidad del libre cambio puede integrarse directamente en la introducción, pues no constituye en absoluto parte esencial del Tema. De desarrollarse por separado, este punto sobre optimalidad del libre comercio de manera mucho más resumida que lo expuesto en el tema (2 minutos, por ejemplo).

     Los apartados sobre Política comercial estratégica y Política de promoción de exterior sí constituyen el esqueleto básico del tema. Cada uno de los mismos recibirá un mínimo de 10 minutos y un máximo de 15 minutos. En opinión de los autores, un reparto equilibrado (12-13 minutos para cada apartado), es aconsejable, pues en el título del Tema no se pide especial referencia a un apartado sobre el otro. Además ello dejaría tiempo (4-6 minutos) para la introducción (con la optimalidad del libre comercio si se desea) y la conclusión.

     En lo concerniente a política comercial estratégica, los autores hemos sido bastante exhaustivos, dado que lo habitual era limitarse a los modelos de tanto de Brander-Spencer como de Eaton y Grossrnan. A la voluntad del opositor queda el grado de profundidad con el que se desarrollan el modelo de Eaton y Kierzowski y él de la industria naciente.

     En lo que se refiere a la parte de promoción exterior es preciso recordar que no es aconsejable el recurso a notables ejemplos de la política comercial española (pese a que ello resultaría muy cómodo al opositor dado que la conoce de la parte A del cuarto ejercicio sobre economía espafiola). Este Tema busca exponer los desarrollos teóricos y prácticos que han influido en la política de promoción (no sólo española sino de otros países).
    }
    
\vspace{1cm}

\modificaciones{
    
    }

\newpage
\MyTOC
\newpage

\mylistofequations
\clearpage

\MyListOfFigures
\clearpage


\pagenumbering{arabic}
\setcounter{page}{1}


\section{Introducción}
Desde la aportación del economista inglés Alfred Marshall a la ciencia económica el estudio teórico de la apertura de los países al exterior se divide en teoría pura del comercio internacional y teoría monetaria del comercio internacional. La teoría monetaria del comercio internacional analiza las economías abiertas desde una perspectiva macroeconómica, esto es, centrándose en conceptos como el tipo de cambio o la balanza de pagos. Por su parte, la teoría pura del comercio internacional, que nos ocupa en este tema, adopta un enfoque microeconómico para explicar los flujos de comercio exterior entre las diferentes naciones en lo que se refiere a volumen, precios, dirección y estructura. Se adopta normalmente el concepto ricardiano de nación como unidades territoriales entre las que existe nula movilidad de factores productivos capital y trabajo. La teoría pura, por tanto, obvia los conceptos de dinero, tipo de cambio y balanza de pagos y se centra en aspectos normativos, como los efectos del comercio sobre el crecimiento y la distribución de la renta o el debate librecambio-proteccionismo. La teoría tradicional (clásica y neoclásica) del comercio internacional asigna superioridad al librecambio frente a cualquier otro tipo de política comercial, pues se produce una especialización eficiente en función de la ventajacomparativa en la tecnología de producción de ciertos bienes (según el principio formulado por David Ricardo y posteriormente refinado por Haberler) o en función de la abundancia factorial (de acuerdo con el modelo de Heckscher-Ohlin-Samuelson). Además, se considera el comercio como un juego de no suma cero donde todos los países pueden salir ganando del intercambio. Los nuevos desarrollos de la teoría del comercio internacional no abogan en su totalidad por un librecambio sin matices como hace la teoría clásica y neoclásica. Existen modelos que otorgan superioridad al librecambio por el acceso a una mejor tecnología y a un mayor elenco de variedades, si bien existen otras teorías que optan por un cierto grado de protección, ya sea de forma temporal o indefinida, el apoyo público a la exportación a través de subvenciones en un contexto estratégico de oligopolio o la apuesta por ciertos sectores en concreto para su promoción exterior. En definitiva, existe un debate sobre la orientación idónea de la política comercial, entendida ésta última como aquella rama de la política económica compuesta de decisiones discrecionales encaminadas a afectar a los flujos de comercio exterior de una economía, ya sea en volumen, precios, dirección o estructura.

El debate de las ideas, en todo caso, parece haber sido ganado por el libre comercio y la mayoría de los economistas han sido manifiestamente librecambistas. Esa optimalidad teórica del librecambio será analizada en el primer punto del tema.


\subsection{Contextualización}


\textcolor{\textbf{OJO}!} -- Resulta interesante mencionar la frase de FRIEDMAN: defensa nacional, sistema jurídico y protección policial. 


\subsubsection{Relevancia}




\subsubsection{Problemática}



\subsection{Estructura de la exposición}
Sin embargo, el librecambismo no ha ganado la batalla de los hechos y ningún país ha aplicado una política pura de libre comercio. Es por ello que el debate librecambio-proteccionismo pemianece vigente y entre ambas posturas intem1edias se sitúa el concepto amplio y ambiguo de política comercial estra1égica, que analizaremos en el segundo punto del tema. De este modo, existirán argumentos que defiendan la protección ante imp011aciones o las subvenciones a la exportación y otro tipo de visiones, que llevan a defender la protección o la promoción de ciertos sectores concretos, aunque con un horizonte de apertura a largo plazo.

Por su parte, el tercer punto del tema estudia la política de promoción exterior como una rama de la política comercial muy utilizada en los países desarrollados que queda habitualmente sin explicar en la mayoría de los modelos de la teoría pura del comercio internacional. La política comercial tiene dos dimensiones muy diferenciadas en función de si los flujos de comercio exterior que se ven afectados son los de importación o los de exportación. La política comercial defensiva comprende todas aquellas acciones dirigidas a reducir las importaciones de una economía, por ejemplo mediante aranceles, restricciones cuantitativas, barreras administrativas, etc.

Por su parte, la política comercial ofensiva abarca todas aquellas medidas destinadas a incrementar las exportaciones de una economía. La política comercial ofensiva tiene, a su vez, dos enfoques: uno tradicional y otro más moderno. 


La política comercial ofensiva tradicional constituye un apoyo público directo a la actividad exportadora, ya sea mediante ventajas de financiación y aseguramiento, subvenciones directas o ventajas fiscales. La política comercial ofensiva más moderna constituye un apoyo genérico directo a la actividad de internacionalización de la empresa. Es en este marco donde se encuadra la denominada política de promoción exterior. En esté tercer punto analizaremos la política de promoción exterior, con una delimitación más exhaustiva de su definición y contenidos. Nos centraremos especialmente en sus objetivos, en su justificación y en los instrumentos al alcance del sector público para lograr esos objetivos.



\clearpage
\section{Política Comercial estratégica}
\puntosclave{
    En los años ochenta y noventa surgieron algunos nuevos argumentos en favor de la intervención gubernamental en el comercio: la teoría de la política comercial estratégica ofrecía razones por las que los países podrían salir ganando si promocionaban determinadas industrias. En los noventa surgió una nueva crítica a la globalización, centrada en los efectos que tenía la globalización sobre los trabajadores de los países en desarrollo. Las posibles acciones sobre el cambio climático han planteado algunas cuestiones importantes sobre el comercio, lo que incluye la deseabilidad y la legalidad de los aranceles sobre el carbono. 
    
    Los argumentos de la política comercial activista descansan sobre dos ideas: (i) el argumento de que los gobiernos deberían promover aquellas industrias que tienen externalidades tecnológicas; y, (ii) es el análisis de BRANDER-SPENCER, que sugiere que la intervención estratégica puede permitir a los países capturar beneficios extraordinarios (que representa un mayor alejamiento de los argumentos habituales sobre los fallos del mercado). Estos argumentos son convincentes en la teoría; sin embargo, muchos economistas consideran que son excesivamente sutiles y exigen demasiada información para resultar útiles en la práctica.

    Se puede analizar la política comercial en los países en desarrollo mediante los mismos instrumentos analíticos que se utilizan para los países avanzados. Sin embargo, la política comercial de los países en desarrollo se ocupa de dos objetivos: promoción de la industrialización y resolución de los desequilibrios en el desarrol1o de la economía nacional.

    La política gubernamental para promover la industrialización ha sido justificada a menudo mediante el argumento de la industria naciente, que sostiene que las industrias nuevas necesitan un periodo de protección frente a la competencia de los competidores establecidos en otros países. 
    
    Al utilizar el argumento de la industria naciente como justificación, muchos países en desarrollo han seguido políticas de industrialización mediante sustitución de importaciones, con las que han creado industrias nacionales bajo la protección de aranceles y cuotas de importación. Aunque estas políticas han tenido éxito en la promoción de las manufacturas, han quedado lejos de proporcionar los beneficios esperados en cuanto a crecimiento económicoy nivel de vida.\footnote{%
        Muchos economistas son hoy severos críticos de los resultados de la sustitución de importaciones, con el argumento de que ha fomentado una producción ineficiente con costes elevados.
    } 
    
    A partir de 1985 muchos países en desarrollo, insatisfechos con los resultados de las políticas de sustitución de importaciones, redujeron en gran medida las tasas de protección de las manufacturas. Con ello, el comercio de esos países en desarrollo creció rápidamente, al igual que la proporción de bienes manufacturados sóbre las exportaciones. Sin embargo, en lo que respecta al crecimiento económico, los resultados de este cambio de política han sido dispares.
}

La política comercial estratégica hace referencia a la actuación de la autoridad fiscal mediante subsidios o impuestos a la exportación o importación con el fin de influir en el resultado de la interacción estratégica entre empresas a favor de la empresa nacional. 

Por tanto exige, como su propio nombre indica, que la estrategia efectuada por la empresa sin la intervención del Estado genere un resultado subóptimo, bien sea en perjuicio de la empresa o en perjuicio del bienestar agregado.\footnote{%
    En situaciones de competencia perfecta la solución óptima es dejar actuar al mercado ya que, con ello, la interacción entre consumidores y empresas en los mercados conduce a un óptimo global en el que el precio de los bienes se iguala a su coste marginal y las empresas obtienen un beneficio económico igual a cero. Sin embargo, en presencia de competencia imperfecta pueden existir beneficios extraordinarios y el precio del bien no iguala al coste marginal, con lo que el nivel de producción es subóptimo y se produce una pérdida de bienestar respecto a la competencia perfecta. De esta idea arranca precisamente la justificación económica de la aplicación de políticas comerciales estratégicas. Efectivamente, los modelos que vamos a ver a continuación demuestran que la aplicación de un impuesto o una subvención a importaciones o exportaciones tiene un efecto neto positivo sobre el bienestar de los ciudadanos del país que lo aplica y de los del mercado que se ve afectado por dicha subvención, si bien el efecto sobre el bienestar global es mucho más discutible. A lo largo de todo este apartado hablaremos de subvenciones e impuestos a la exportación, pero nuestras conclusiones serian similares si en lugar de subvencionar o gravar la exportación se decidiera subvencionar o gravar la producción en el supuesto de que las empresas no venden en su país y sólo lo hacen en un país tercero.
} Esta condición lleva a que el análisis de este tipo de intervención se realice sobre el marco de la competencia oligopolística: un conjunto de empresas (nacionales y extranjeras) que dominan en el mercado del producto (homogeneo o heterogeneo) e internalizan las decisiones del resto de empresas en la toma de las propias, bien sea la ejecución de un proyecto de inversión, el nivel de producción o los precios. 

Desde este punto de vista existen diferentes clasificaciones. No obstante, para aterrizar el resultado a la experiencia real veremos dos situaciones fundamentales en las que la intervención puede ser justificada por motivos estratégicos.\footnote{%
    El uso de distintos modelos deriva del hecho de que no existe un modelo universalmente aceptado de oligopolio y, menos aún, en el contexto de la Política Comercial. Por ello, hemos de contemplar varios posibles modelos que varían sustancialmente en sus supuestos de partida y que por ello dan lugar a conclusiones diferentes cuando se aplican a la política comercial estratégica.
}

\subsection{Promoción: Ventaja estratégica}
En primer lugar encontramos la situación en la que a través de la intervención pública se busca proporcionar una ventaja estratégica a la empresa nacional. Se trata del marco convencional de estudio iniciado por los economistas SPENCER y BRANDER.

Estos apuntan a que, en algunas industrias, solo un número reducido de empresas compite realmente entre sí. Debido al pequeño número de compañías implicadas, no son aplicables los supuestos de competencia perfecta: normalmente habrá beneficios extraordinarios. Es decir, las empresas obtendrán beneficios por encima de lo que las inversiones del mismo riesgo podrían conseguir en cualquier sector de la economía. Ello conducirá a una competencia internacional para apropiarse de esos beneficios. SPENCER y BRANDER observaron que, en este caso, es posible, en principio, que un gobierno altere las reglas del juego y traslade estos rendimientos excesivos del extranjero a las empresas nacionales. En el caso más sencillo, un subsidio a las empresas nacionales, que desalienta la inversión y la producción de las competidoras extranjeras, puede aumentar, los beneficios de las empresas nacionales en una cantidad superior al subsidio. Si se omiten los efectos sobre los consumidores (por ejemplo, cuando las empresas venden solamente en mercados extranjeros), esta captura de beneficios de los competidores extranjeros significaría que el subsidio incrementa la renta nacional a expensas de otros países. 

Veamos dos ejemplos.

\subsubsection{Promoción: Costes hundidos}
Supongamos que existen dos empresas compitiendo, cada una de un país distinto. Llamaremos a las empresas Boeing y Airbus, y a los países Estados Unidos y Europa. Sin entrar en los detalles propios del sector aeroespacial, si podemos intuir algunas características especialmente relevantes: (i) se trata de un sector con altos costes de I+D; (ii) genera, normalmente, muchos efectos desbordamiento positivos (spillovers); y, (iii) puede considerarse que se trata de un sector que reune las características básicas para considerarse estratégico (y, por tanto, adecuado preservar). Desarrollar una nueva gama de aviones comerciales toma mucho tiempo (alrededor de unos 10 años de proyecto), y los beneficios, aunque cuantiosos, no están completamente garantizados. Esto se debe, en mayor parte, porque en este sector no hay pan para todos, normalmente el ganador se lleva la mayoría de los contratos de producción. 

Supongamos que ambas empresas desean desarrollar una nueva gama de aviones comerciales que las dos empresas son capaces de producir. Para mayor sencillez, consideraremos que cada empresa puede tomar solo la decisión de sí, o no. A través de la represetación normal del juego podemos definir la matriz de pagos de la siguiente forma:

\imagenfit{Airbus_BOEING_juego.png}{Competencia de dos empresas}{}

Como vemos, muestra la dependencia de los beneficios obtenidos por las dos empresas de sus decisiones. Cada fila corresponde a una decisión particular de Boeing; cada columna, a una decisión de Airbus. En cada cuadro hay dos entradas: la entrada inferior izquierda representa los beneficios de Boeing, mientras que la superior derecha representa los beneficios de Airbus. Cada empresa sola podría tener beneficios con el proyecto, pero si las dos empresas los producen, ambas tendrán pérdidas. ¿Qué empresa obtendrá realmente beneficios? Ello depende de quién tome la iniciativa. 

Supongamos que Boeing (empresa lider) puede tener una pequeña ventaja inicial, y decide producir antes de que Airbus (empresa seguidora) pueda decidirlo. En este escenario, Airbus no tendrá incentivos para entrar y Boeing acabará obteniendo todas las rentas del proyecto (en el juego de STACKELBERG, maximiza sus beneficios a expensas de la seguidora). El resultado se situará en la parte superior derecha de la tabla y será Boeing quien obtiene los beneficios. 

Ahora entra en juego la observación de BRANDER-SPENCER: el gobierno europeo puede revertir esta situación. Supongamos que dicho gobierno europeo decide pagar a su empresa un subsidio de 25 si entra en el negocio. Como consecuencia se producirá un cambio en el cuadro de resultados:

\imagenfit{Airbus_BOEING_juego_subsidio.png}{}{}

Gracias al subsidio, es rentable que Airbus inicie el proyecto, con independencia de lo que haga Boeing. ¿Por qué? Boeing sabe que, de cualquier manera, ha de competir con Airbus y, por tanto, perderá dinero si elige producir. Por consiguiente, será Boeing quien se verá disuadida de entrar en el negocio. Por tanto, el subsidio del gobierno ha modificado la ventaja del liderazgo inicial, que suponíamos que tenía Boeing, y la ha otorgado a Airbus. El resultado final es que el subsidio incrementa los beneficios en una cantidad mayor que el propio subsidio, debido a su efecto disuasorio sobre la competencia extranjera. El subsidio tiene estos efectos porque \textbf{crea una ventaja para Airbus comparable con la ventaja estratégica que habría tenido si ella, y no Boeing, hubiera alcanzado un liderazgo inicial en la industria} (por eso lo más común es presentarla en un juego tipo STACKELBERG).

\paragraph{Limitaciones: ¿Cuánto sabemos realmente?}
La justificación estratégica en favor de la polftica comercial, aunque ha atraído mucho interés, también ha recibido muchas críticas. Los críticos consideran que, para hacer un uso práctico de la teoría, se requeriría más información de la que probablemente se dispone, que tales políticas corren el riesgo de provocar represalias extranjeras y que, en cualquier caso, los aspectos políticos de la política comercial y de la industrial impedirían el uso de herramientas analíticas tan sutiles. 

El problema de la información insuficiente tiene dos vertientes. La primera es que, incluso si se considera a una industria de forma aislada, puede ser difícil rellenar las entradas en una tabla con alguna confianza. Por otra parte, si el gobierno lo hace mal, una política de subsidio puede terminar por convertirse en un costoso error. Supongamos que Boeing ofrece alguna ventaja inherente (una mejor tecnología, menores costes de ajuste, etc.), por lo que, aunque Airbus entre en el negocio, será beneficioso para Boeing producir. 

\imagenfit{Airbus_BOEING_juego_VentajaIntrinseca.png}{Competencia entre dos empresas: Ventaja intrínseca de BOEING}{Elaboración propia}

Como vemos, si Boeing verdaderamente tiene una ventaja intrínseca, existen dos Equilibrio de NASH y no puede garantizarse con seguridad que Boeing no vaya a producir en caso de que Airbus produzca. No obstante, Airbus no puede producir con beneficios si Boeing entra, por lo que, sin un subsidio, el resultado será que Boeing produce y Airbus no. Debido a la ventaja subyacente de Boeing, el subsidio no actuará como un elemento disuasorio para Boeing, y los beneficios de Airbus serán inferiores a la cuantía del subsidio; en definitiva, la política habrá sido un caro error. La cuestión es que, aun cuando los dos casos son muy parecidos, en uno el subsidio podría ser una buena idea, mientras que en el otro no parece tener ningún sentido. Es como si la deseabilidad de las políticas comerciales estratégicas dependiera de una lectura exacta de la situación. Ello conduce a algunos economistas a preguntarse si hay siempre suficiente información para utilizar la teoría de manera eficaz. La necesidad de información se complica por el hecho de que no es posible considerar las industrias de forma aislada. Si una industria es subsidiada, detraerá recursos de otras y producirá un aumento de sus costes. Así, incluso una política que tenga éxito en la consecución de una ventaja estratégica en una industria para las empresas de los Estados Unidos tenderá a provocar desventajas estratégicas en otras. 

Para saber si la política está justificada, el gobierno de los Estados Unidos debe ponderar esos efectos compensatorios. Aunque los poderes públicos tengan un conocimiento exacto de una industria, este no es suficiente: necesitan un conocimiento igualmente preciso de todas las industrias con las que esta industria compite por los recursos. Aun cuando una política comercial estratégica propuesta pueda superar estas críticas, deberá hacer frente todavía al problema de la represalia exterior, esencialmente el mismo problema al que nos enfrentábamos cuando considerábamos la aplicación de un arancel para mejorar la relación de intercambio.\footnote{%
    Las políticas estratégicas son políticas para empobrecer al vecino, que aumentan nuestro bienestar a expensas de otros países. Por tanto, estas políticas corren el riesgo de provocar una guerra comercial que perjudique a todo el mundo. Pocos economistas serian favorables a que los Estados Unidos iniciaran este tipo de políticas. Más bien, lo máximo que se defiende normalmente es que la nación estadounidense se prepare para tomar represalias cuando otros países apliquen politicas estratégicas de forma agresiva.
}

\subsubsection{Restricción: Excesivamente \textit{competitivo}}
Otra situación en la que nos podemos encontrar una desventaja estratégica que perjudica la empresa nacional y exporta rentas a la empresa extranjera es cuando existen conjeturas mal formadas.\footnote{%
    Los modelos de variaciones conjeturales con profundamente complejos desde un punto de vista conceptual. Se ha tendido, sistemáticamente, a su representación gráfica, lo cual constituye (en opinión del autor que ha sufrida esto) una profunda desventaja a la hora que enfrentamos la verdadera naturaleza de este fenómeno. Para conocer tanto la derivación analítica de la conjetura como los debates académicos al respecto: [\authorfont{Ver Tema 3.A.19}].
} Una de las implicaciones más importantes (sino la más) de las conjeturas, en la competencia oligopolística, es su impacto sobre la estructura de mercado. 

Es decir, la creencia de que una empresa sobrerreacionará ante desviaciones propias caracterizará un mercado con beneficios extraordinarios, menor producto y mayores precios. Por el contrario, la creencia de que una empresa será pasiva ante desviaciones propias caracterizará un mercado mucho más cercano próximo a una estructura competitiva (con beneficios extraordinarios), mayor producto y menores precios. Las razones por las que una empresa forma una u otra conjetura y, en consecuencia, el evidente impacto que esto tiene sobre sus rentas y la estructura económica son muy dispares. Ahora bien, en Comercio Internacional esto tiene una importante consecuencia y puede justificar otro tipo de intervención por razones estratégicas. 

Supongamos que un mercado está formado por un pequeño número de empresas distribuidas entre diferentes territorios. Nuestra empresa doméstica conjetura que el resto de empresas del sector serán pasivas ante los intentos de esta por apropiarse de gran parte del mercado. En esta situación, es de esperar que nuestra empresa doméstica ampliaría su producción, lo que provocaría una disminución del precio internacional en dicho mercado. A priori esto puede ser un buen resultado desde una óptica nacional pero, en realidad, no lo es tanto. En primer lugar, los beneficios de este incremento de la producción y consecuentemente bajada del precio se distribuyen más allá de nuestras fronteras. En segundo lugar, los beneficios de la empresa pueden ser menores de los que alcanzaría en una situación de colusión más tácita. Más aún si consideramos, por ejmplo, que dicho mercado es además considerado estratégico desde un punto de vista nacional (por ejemplo, supongamos que se trata de la exportación de tierras raras, abundantes en nuestra nación, pero finitas). ¿Qué debemos hacer?

En este escenario, la Política Comercial Estratégica resulta de gran utilidad. En puridad, nuestra empresa no está en una mejor posición estratégica, puesto que subestima la reacción de sus competidores y, en consecuencia, no prevée correctamente la reacción dinámica de estos. Por ello, lo óptimo desde un punto de vista social sería recurrir a los impuestos a la exportación. Esto es, deberíamos reducir el impacto de esta mala conducta estratégica forzando a la empresa doméstica a bajar su producción. Lo que, de hecho, sucede en el mercado de las tierras raras, donde los Gobiernos nacionales exigen estrictas cuotas de exportación y aprovisionamiento. 

\subsection{Protección: ------}
Vistos dos de los argumentos más comúnmente utilizados para justificar la intervención para la promoción estratégica de una empresa doméstica, veamos otra de las justificaciones más importantes desde una persepectiva histórica: la protección estratégica de la industria naciente.


\subsubsection{Polícias de Industrialización por Sustitución de Importaciones}

\textcolor{red}{\textbf{OJO} -> Aquí el desarrollo teórico de la Industralización por Sustitución de Importaciones}

A medida que un número creciente de países obtuvo su independencia de las potencias coloniales en el período inmediatamente posterior al final de la Segunda Guerra Mundial, se desarrolló entre economistas y responsables de política económica la opinión generalizada de que la mejor manera de que estos países se desarrollaran más rápidamente era estimular la industrialización mediante la adopción de políticas de sustitución de importaciones. Para estos, la industrialización parecía ofrecer la posibilidad de lograr un crecimiento más rápido, mayores niveles de renta per cápita y el poder económico y militar necesario para la seguridad nacional. Una forma económicamente sensata de lograr la industrialización parecía ser restringir las importaciones de bienes manufacturados para los cuales ya existía demanda interna, con el fin tanto de desplazar dicha demanda hacia los productores nacionales como de permitir el uso de los ingresos por exportaciones de productos primarios del país para importar los bienes de capital necesarios para la industrialización.\footnote{%
    También parecía haber varios ejemplos en los que altos niveles de protección a las importaciones en los siglos XIX y XX habían contribuido positivamente a la industrialización. Aunque Gran Bretaña había adoptado una política de libre comercio durante su período de rápido crecimiento en el siglo XIX, Estados Unidos parecía haberse industrializado y prosperado imponiendo elevados aranceles a las manufacturas durante gran parte de la segunda mitad del siglo XIX. Alemania y Francia también adoptaron políticas proteccionistas durante este período, al igual que Japón después de 1900. El notable grado de industrialización alcanzado por Unión Soviética en las décadas de 1920 y 1930 y por China después de 1949 mediante la adopción de políticas orientadas hacia el interior fueron ejemplos históricos adicionales que impresionaron a los líderes de las naciones recién independientes.
}

\begin{specialblock}{El argumento de la industria naciente}
    El llamado argumento de la industria naciente, formulado por primera vez en 1791 por Alexander Hamilton, desarrollado posteriormente por Friedrich List (1856), y aceptado por muchos economistas clásicos y neoclásicos como la principal excepción teóricamente válida al caso del libre comercio mundial, proporcionó apoyo económico a las políticas de sustitución de importaciones. John Stuart Mill, quien formalizó por primera vez el argumento en términos económicos, sostenía que se necesita tiempo para que los nuevos productores en un país se “eduquen hasta el nivel de aquellos para quienes los procesos son tradicionales” y, por tanto, para que sus costes unitarios disminuyan. El argumento de la industria naciente sostiene que, durante el período temporal en el que los costes internos en una industria están por encima del precio de importación del producto, un arancel es un método socialmente deseable de financiar la inversión en recursos humanos necesaria para competir con éxito con los productores extranjeros. 
    
    Poco después de la Segunda Guerra Mundial, Raúl Prebisch (1950), quien fue Secretario General de la Comisión Económica para América Latina de las Naciones Unidas y posteriormente organizó y fue Secretario General de la Conferencia de las Naciones Unidas sobre Comercio y Desarrollo (UNCTAD), entre otros, argumentó que el argumento de la industria naciente era aplicable a todo el sector manufacturero y no solo a una industria específica. También afirmó que una tendencia secular a la baja en los precios de los productos primarios (las exportaciones de los países menos desarrollados) en relación con los precios de los bienes manufacturados (las exportaciones de los países desarrollados), junto con la baja elasticidad de la demanda de productos primarios, hacía poco atractiva la expansión de la producción de estos. Centrarse en producir bienes manufacturados intensivos en trabajo, por ejemplo, ropa, con fines de exportación tampoco resultaba atractivo para la mayoría de los países menos desarrollados en ese momento, debido a la creencia de que una estructura industrial equilibrada, como la existente en la mayoría de los países desarrollados, era necesaria para alcanzar altos niveles de renta per cápita y, además, porque persistían altos niveles de aranceles y otras barreras a la importación en los países desarrollados para la mayoría de estos bienes. 
    
    Aunque la mayoría de los líderes económicos de los países menos desarrollados veían favorablemente la estrategia de sustitución de importaciones, a menudo se encontraron empujados hacia dicha política de manera algo inadvertida. Debido a la escasez de bienes que estos países sufrieron durante la Segunda Guerra Mundial y a los planes de expansión económica de sus nuevos líderes, existía una enorme demanda tanto de bienes de capital como de bienes de consumo. Esto implicaba que sus reservas de divisas existentes se agotaban rápidamente, siendo los ingresos corrientes por exportaciones insuficientes para cubrir la brecha entre demanda y oferta a los tipos de cambio vigentes. En consecuencia, la mayoría de estos países se vio obligada a imponer controles de divisas y de importaciones para conservar sus ingresos por exportaciones y establecer un sistema de racionamiento de las divisas disponibles, con el fin de asegurar que bienes de consumo esenciales como alimentos y medicinas, insumos intermedios clave como combustibles, y bienes de capital esenciales pudieran importarse en cantidades suficientes para evitar graves disturbios políticos y, al mismo tiempo, permitir la consecución de sus objetivos de desarrollo. Una consecuencia fue que se establecieron niveles muy altos de protección implícita sobre los llamados bienes manufacturados “no esenciales”.
\end{specialblock}

Las políticas de sustitución de importaciones funcionaron bastante bien inicialmente. Los altos precios de los bienes no esenciales importados desplazaron la demanda interna de estos bienes desde los productores extranjeros hacia los locales, con el resultado de aumentos significativos en la producción de manufacturas simples, ya que los gobiernos proporcionaban a los productores nacionales las divisas necesarias para importar insumos intermedios clave y bienes de capital.\footnote{%
    Muchas actividades manufactureras consistían en gran medida en ensamblar componentes producidos en el extranjero, por ejemplo, automóviles. Dado que la producción de la mayoría de estos bienes utilizaba intensivamente el tipo de trabajo relativamente abundante en las naciones recientemente industrializadas, los efectos adversos sobre la eficiencia económica de estos primeros esfuerzos de sustitución de importaciones no fueron suficientes para compensar sus efectos positivos sobre el crecimiento. Además, en este período inicial, las monedas nacionales sobrevaluadas resultantes de los estrictos controles cambiarios y de políticas productivas expansivas no parecían reducir significativamente los ingresos por exportaciones de productos primarios y mantenían relativamente bajos los precios de importación de bienes de capital e insumos intermedios necesarios.
} No obstante, a medida que las políticas de sustitución de importaciones continuaron y varios países en desarrollo extendieron estas políticas para abarcar cada vez más insumos intermedios y bienes de capital, las desventajas de este enfoque se hicieron cada vez más evidentes.\footnote{%
    En particular, las dificultades impuestas al sector exportador comenzaron a tener efectos adversos sobre el crecimiento. Una moneda sobrevaluada implicaba que el número de unidades de divisas recibidas por los exportadores permanecía bajo, mientras que, al mismo tiempo, estos productores se veían obligados a adquirir cada vez más insumos intermedios y bienes de capital a nivel nacional a precios elevados. La consiguiente compresión de los márgenes de beneficio los obligó a reducir la producción destinada a la exportación. Los mayores requisitos de cualificación y tecnología para producir bienes intermedios y de capital más complejos, junto con la falta de mercados internos amplios necesarios para alcanzar niveles eficientes de producción, también empeoraron las perspectivas de beneficios para los productores nacionales. Al mismo tiempo, las políticas expansivas agresivas de los gobiernos y de las empresas privadas alimentaron mayores presiones inflacionarias, con el resultado de que los déficits presupuestarios públicos y los déficits en la balanza de pagos se volvieron comunes. Las crisis presupuestarias y de balanza de pagos resultantes a menudo se enfrentaron con controles aún más estrictos sobre los tipos de cambio y las importaciones, así como con una mayor intervención del gobierno en la economía.
} El resultado neto fue, en general, una desaceleración de la tasa de crecimiento en comparación con el período inicial de sustitución de importaciones.

Dado el amplio consenso actual entre los economistas de que la estrategia de sustitución de importaciones no funcionó bien para la mayoría de los países en desarrollo, surge una cuestión importante: por qué tantos economistas se equivocaron en sus predicciones de que tal enfoque tendría éxito en elevar las tasas de crecimiento a largo plazo de estos países. ¿Qué falló en nuestro razonamiento analítico? En mi opinión, dos errores fueron una aceptación acrítica del argumento de la industria naciente y una falta de consideración de las consecuencias macroeconómicas de dicha política cuando se aplica a todo el sector manufacturero.

Considérese el argumento expuesto anteriormente de que los nuevos productores necesitan ser protegidos durante un período temporal para poder adquirir la experiencia y las habilidades productivas que los harán tan eficientes como sus competidores extranjeros establecidos. Como señaló James Meade (1955) hace muchos años, la existencia, durante el período inicial de producción, de costes superiores a los de los competidores extranjeros no constituye, por sí sola, una razón suficiente para justificar la protección arancelaria desde el punto de vista de la eficiencia económica. Si los costes unitarios en una industria son lo suficientemente más bajos después del período de aprendizaje como para generar un excedente descontado de ingresos sobre costes (indicando así una ventaja comparativa para el país en la producción del bien), debería ser posible para las empresas obtener fondos suficientes en el mercado de capitales para cubrir su exceso inicial de gastos sobre ingresos. Estas circunstancias no son diferentes de aquellas en las que las empresas acuden al mercado de capitales para financiar el exceso de gastos sobre ingresos en las primeras etapas de producción debido a la necesidad de adquirir unidades indivisibles de capital físico. Las imperfecciones en los mercados de capitales pueden impedir el acceso a estos, pero la existencia de tales imperfecciones constituye un caso distinto para la intervención gubernamental del que plantea el argumento de la industria naciente.

Como también señaló Meade (1955), el argumento clave sobre el que debe basarse el caso de la industria naciente se refiere a las externalidades tecnológicas asociadas al proceso de aprendizaje. Por ejemplo, considérese la adquisición del conocimiento sobre técnicas de producción locales necesario para competir eficazmente con productores extranjeros. Un empresario que incurre en estos costes de descubrimiento de la mejor forma de producir un bien se enfrenta al problema de que esta información puede volverse libremente accesible para otros productores locales potenciales, quienes pueden utilizarla simultáneamente sin incurrir en los costes completos de adquisición del conocimiento. La competencia de estos otros productores podría entonces elevar los precios de los factores o reducir el precio del producto hasta niveles en los que la empresa inicial no pueda recuperar sus costes de obtención de ese conocimiento. Anticipando este resultado, las empresas se verán desincentivadas a asumir los costes iniciales de adquisición de conocimiento.

La imposición de un arancel protector temporal no garantiza, sin embargo, que los empresarios individuales realicen inversiones adicionales en adquisición de conocimiento. Un impuesto a las importaciones eleva el precio interno de un producto y, desde el punto de vista de la industria en su conjunto, hace más rentables algunas inversiones en conocimiento. Pero los productores individuales siguen enfrentándose al mismo problema de externalidad que antes: otras empresas copiarán, con poco coste, cualquier nuevo conocimiento técnico descubierto y reducirán el precio del producto hasta un nivel en el que la empresa inicial no pueda recuperar sus costes. Si existiera siempre un desfase temporal tecnológicamente determinado entre la introducción de una nueva técnica de producción más barata y el ajuste de precios causado por la entrada de empresas imitadoras, un arancel haría más rentable la inversión en conocimiento para la empresa individual. Sin embargo, la rapidez con la que las empresas responden a las oportunidades de mercado depende a su vez del nivel de beneficios esperados. Un arancel incentivará a las empresas a incurrir más rápidamente en los costes de adquisición del conocimiento descubierto por otras y a entrar en la producción con mayor rapidez y a niveles elevados de producción. Lo que se requiere, en realidad, es un subsidio a los primeros entrantes en la industria con el fin de descubrir mejores técnicas de producción.

Hasta el período posterior a la Segunda Guerra Mundial, cuando algunos economistas comenzaron a extender el argumento de la industria naciente a todo el sector manufacturero, este argumento se formulaba generalmente en términos de equilibrio parcial. Se centraba en una sola industria y se asumía que la protección temporal a las importaciones no tenía efectos apreciables sobre variables macroeconómicas como los tipos de cambio, las exportaciones e importaciones agregadas, o las políticas monetarias y fiscales. Los primeros defensores de la protección agresiva de amplios segmentos del sector manufacturero no apreciaron plenamente las implicaciones de sus propuestas sobre estas variables macroeconómicas. No tuvieron suficientemente en cuenta, por ejemplo, los efectos adversos de la sustitución de importaciones sobre las exportaciones agregadas y, por tanto, sobre los ingresos en divisas necesarios para importar bienes de capital e insumos intermedios esenciales. Tampoco comprendieron hasta qué punto las políticas destinadas a conservar divisas mediante la restricción de importaciones de bienes de consumo de lujo harían que la producción interna de estos bienes fuera la opción más atractiva para los empresarios, sesgando así la estructura productiva en una dirección no deseada. Asimismo, subestimaron las presiones presupuestarias e inflacionarias generadas por las políticas de desarrollo. De hecho, fueron las crisis macroeconómicas asociadas a déficits de importación insostenibles, déficits fiscales incontrolables e inflación desbordada las que finalmente llevaron a la mayoría de los países a abandonar la sustitución de importaciones, más que la constatación de los problemas de asignación de recursos derivados de estas políticas.


\textcolor{red}{\textbf{OJO} -> A partir de aquí Manual de KRUGMAN, OBSTFELD y MELITZ} (Juntas a ser posible)

Durante unos 30 años tras la Segunda Guerra Mundial, en muchos países en desarrollo influyó con fuerza la creencia de que la clave para el desarrollo económico era la creación de un sector industrial potente, y que la mejor forma para crear ese sector industrial era mediante la protección de los fabricantes nacionales frente a la competencia internacional. Hablamos en particular de lo que se conoce comúnmente como la industrialización mediante la sustitución de importaciones.

Este camino a seguir se sustenta sobre el argumento de la industria naciente. Según este, los países en desarrollo tienen una ventaja comparativa potencial en la producción de manufacturas, pero las nuevas industrias manufactureras en los países en desarrollo no pueden competir, inicialmente, con las manufacturas establecidas anteriormente en los países desarrollados. Por eso, para permitir que las manufacturas tomen impulso, los gobiernos deberían apoyar temporalmente las nuevas industrias, hasta que tengan un tamaño suficiente para enfrentarse a la competencia internacional. Así pues, tiene sentido, según este argumento, utilizar aranceles o cuotas de importación como medidas transitorias para permitir el inicio de la industrialización.\footnote{%
    Es un hecho histórico que algunas de las mayores economías de mercado del mundo iniciaron su industrialización sustentadas en barreras comerciales: los Estados Unidos tuvieron elevadas tasas arancelarias en tas manufacturas durante el siglo XIX, mientras que Japón impuso amplios controles de importación hasta los afios setenta.
} 

El argumento de la industria naciente parece muy posible y ha persuadido, de hecho, a muchos gobiernos. Sin embargo, los economistas han encontrado muchos peligros en este argumento y han sugerido que debe ser utilizado con cautela. En primer lugar, no siempre es buena idea intentar avanzar hoy hacia industrias que tendrán una ventaja comparativa en el futuro. Supongamos que un país que es abundante en trabajo está en un proceso de acumulación de capital. Cuando acumule suficiente capital tendrá ventaja comparativa en industrias intensivas en capital. Sin embargo, esto no significa que deba intentar desarrollar esas industrias inmediatamente.\footnote{%
    En los años ochenta, por ejemplo, Corea del Sur se convirtió en exportador de automóviles; probablemente no hubiera sido una buena idea que esta nación desarrollara su industria automovilística en los años sesenta, cuando el capital y la mano de obra cualificada eran todavía muy escasos.
} En segundo lugar, la protección de la producción de manufacturas no es buena, a menos que la propia protección ayude a hacer que la industria se haga competitiva.\footnote{%
    Por ejemplo, Pakistán y la India han protegido sus sectores manufactureros durante décadas, y solo recientemente han comenzado a realizar exportaciones de bienes manufacturados de cierta importancia. Los bienes que exportan, sin embargo, son manufacturas ligeras como textiles, no las manufacturas pesadas que protegen; así pues, se podría pensar que habrían desarrollado sus exportaciones manufactureras, aunque no hubieran protegido nunca su industria.
} 

\paragraph{Justificación a la intervención: Fallos de mercado}
De forma más general, el hecho de que es costoso y requiere tiempo desarrollar una industria no es un argumento para la intervención pública, a no ser que haya fallos en el mercado nacional.\footnote{%
    Algunos economistas han llamado la atención sobre el caso de la industria seudonaciente en el que la industria es protegida inicialmente y, entonces, se hace competitiva por razones que no tienen nada que ver con la protección. En este caso, la protección de la industria naciente finaliza con un éxito aparente, pero en realidad puede haber tenido un coste neto para la economía.
} Es decir, el argumento para proteger una industria en su nacimiento debe estar relacionado con fallos que impidan a los mercados privados desarrollar la industria con la rapidez que debieran; ya que, si se supone que una industria es capaz de permitir ganancias suficientemente elevadas para el capital, el trabajo y otros factores productivos como para que merezca la pena su desarrollo, entonces, ¿por qué no desarrollan dicha industria los inversores privados sin ayuda de la administración?

Estrechamente relacionaco con esto es el primer argumento a través del cual muchos economistas han tratado de justificar dicha protección: las imperfecciones del mercado de capitales.

Si un país en desarrollo no tiene un conjunto de instituciones financieras (como bancos y mercados de valores eficientes) que permitan que el ahorro de los sectores tradicionales (como la agricultura) sea utilizado para financiar la inversión en sectores nuevos (como la producción de manufacturas), entonces el crecimiento de nuevas industrias se verá restringido por la capacidad de las empresas en dichas industrias para obtener beneficios a corto plazo. Así, los bajos beneficios iniciales serán un obstáculo para la inversión, a pesar de que los beneficios a largo plazo de dicha inversión sean altos.\footnote{%
    También se argumenta que los inversores privados en una industria tienen en cuenta solamente los beneficios actuales y olvidan las perspectivas futuras, aunque este razonamiento no es consistente con la conducta del mercado. Al menos en los países avanzados, donde los inversores impulsan a menudo proyectos cuyos beneficios son inciertos y lejanos en el futuro. Considérese, por ejemplo, la industria de la biotecnología en los Estados Unidos, que atrajo cientos de millones de dólares de capital años antes de que se produjera una sola venta comercial.
}

En este escenario, la política óptima es crear un mercado de capitales mejor; no obstante, la protección de las nuevas industrias, que aumentarían sus beneficios y, de este modo, lograrían un crecimiento más rápido, puede estar justificada como opción política de segundo óptimo.

Otro argumento al que se alude con mucha frecuencia es el de la apropiabilidad. Este argumento puede adoptar múltiples formas, todas las cuales comparten la idea de que las empresas en una industria nueva generan beneficios sociales por los que no obtienen compensación. 

Por ejemplo, las empresas que acceden primero a una industria pueden haber incurrido en los costes de establecimiento de adaptación tecnológica a las circunstancias locales o de apertura de nuevos mercados. Si otras empresas las siguen sin incurrir en estos costes de establecimiento, las pioneras no podrán reclamar los beneficios de dicha inversión. Así pues, las empresas pioneras pueden, además de realizar la producción física, crear beneficios intangibles (como conocimientos y nuevos mercados) sobre los que no pueden establecer derechos de propiedad. En algunos casos, los beneficios sociales de la creación de una industria nueva excederán a los costes; a pesar de todo, precisamente por el problema de la apropiabilidad, la iniciativa privada no querrá entrar. 

En este escenario, la respuesta óptima consistiría en compensar a las empresas por sus contribuciones intangibles. Sin embargo, cuando esto no es posible, hay un segundo óptimo para fomentar la entrada en nuevas industrias, mediante la imposición de aranceles u otras políticas comerciales.

\paragraph{Implicaciones de Política Económica}
En definitiva, el argumento de la industria naciente es considerado válido para muchos economistas si y solo si se puede justificar por la existencia de un fallo del mercado a favor de la intervención. Ambos argumentos son ejemplos claros en que los fallos del mercado justificarían interferir en el libre comercio. 

La diferencia es que, en este caso, los argumentos se aplican concretamente a industrias nuevas, no a cualquier industria. No obstante, aún persisten problemas generales del enfoque de los fallos de mercado. En la práctica resulta dificil evaluar qué industrias justifican un trato especial, y el riesgo de que una política que intente promover el desarrollo acabe por caer presa de determinados intereses. Se han dado casos numerosos de industrias nacientes que nunca crecieron y mantuvieron su dependencia de la protección.

De hecho, este punto de vista fue rápidamente abandonado por la escasa evidencia en su favor y, en particular, tras el rápido crecimiento económico observado en varias ecémomías asiáticas. Estas economías se han industrializado no por sustitución de importaciones, sino por las exportaciones de bienes mánufacturados. Se caracterizan por altas ratios de comercio con respecto a la renta nacional y por tasas de crecimiento extremadamente elevadas. Aunque, cabe destacar, que las razones del éxito de estas economías son muy controvertidas y existe un encendido debate en torno al papel que ha desempeñado la liberalización comercial.

\subsubsection{Políticas de Industrialización orientada a la exportación}



\paragraph{La planificación estratégica: CHINA}


\paragraph{Implicaciones de Política Económica}








\clearpage
\section{Política de Promoción Exterior}
\puntosclave{
    
}

\textcolor{red}{\textbf{OJO!}} -- Aquí puede ser interesante poner los dos ejemplos: \textbf{Japón y semiconductores} para ilustrar barreras sociales, culturales, etc. y, por otro lado, \textbf{POLITICO BeBartlet} para ilustrar impedimentos locales relacionados con aspectos culturales y de aptitudes (p.e. PYMES, estructura básica en España).


En lo concerniente a la política de promoción exterior, el recurso a la misma es claramente más habitual en los países desarrollados. Es fundamental por lo tanto entender las razones para el desempeño de esta política por parte del sector público. Por esta razón hemos analizado los fallos de mercado que afectan a la decisión de internacionalización. No obstante, estos fallos de mercado son condición necesaria pero no suficiente para la intervención pública. La racionalidad de la participación del sector público en este campo debe matizarse en función de un análisis coste-beneficio.


La política comercial ofensiva abarca todas aquellas medidas destinadas a incrementar las exportaciones de una economía. La política comercial ofensiva tiene, a su vez, dos enfoques: uno tradicional y otro más moderno. La política comercial ofensiva tradicional constituye un apoyo público directo a la actividad exportadora, ya sea mediante ventajas de financiación y aseguramiento, subvenciones directas o ventajas fiscales. La política comercial ofensiva más moderna constituye un apoyo genérico directo a la actividad de internacionalización de la empresa. Es en este marco donde se encuadra la denominada política de promoción exterior. Encontramos en esta definición dos diferencias entre la política comercial ofensiva tradicional y moderna. La primera diferencia es la condición de apoyo genérico, poniendo un marco de instrumentos a manos de las empresas, pero sin ser un apoyo vinculado directamente al volumende exportación.La segunda diferencia es el hincapié en el concepto de internacionalización, más amplio que él de la política tradicional. Es decir, la política ofensiva tradicional puede generar un éxito ocasional y puntual de una empresa exportadora, pero dicho éxito está ligado directamente a ese apoyo oficial condicionado al aumento de las exportaciones. Por su paite, la política ofensiva moderna busca apoyar a las empresas en una primera fase del proceso de internacionalización para que éstas establezcan lazos sólidos, estables y duraderos y no dependan en el futuro del apoyo públ ico. El hecho de utilizar el concepto de internacionalización nos lleva a recordar que la política de promoción exterior se dedica tanto a la promoción de exportaciones como al fomento de la emisión de Inversión Directa Extranjera. Sin embargo, sus políticas no van a variarsignificativamente en función de la naturaleza exportadora o inversora de la empresa afectada, sino que van a poner instrumentos para que cada empresa se autoseleccione y elija si quiere exportar a un mercado exterior, invertir en el mismo o ambas cosas. Lógicamente, la política comercial ofensiva tradicional no sólo apoya a empresas exportadoras sino también a inversoras, pero las considera como compartimentos estancos y diseñará, por un lado, instrumentos de apoyo público a la exportación (subvenciones) para empresas exportadoras y, por otro lado, instrumentos de apoyo público a la inversión para empresas inversoras. Ello contrasta con el amplio concepto de internacionalización, donde caben empresas que pueden variar su estatus de exportadora a inversora continuamente y que muy a menudo serán ambas cosas. Este tipo de empresas pueden requerir un apoyo genérico a su internacionalización (por ejemplo, mediante información sobre






















































\clearpage
\section{Conclusión} 


\subsection{Recapitulación}


\subsection{Extensiones y relación con otras partes del Temario}



\subsection{Opinión}

\subsubsection{Idea final (Salida o cierra)}

\clearpage
\pagenumbering{gobble}
\MyBibliography



\MyBibItem{DAWID y GATTI}{Chapter 2 - Agent-Based Macroeconomics}{2018}{https://doi.org/10.1016/bs.hescom.2018.02.006}


% ANEXOS
\clearpage
\anexo{\textbf{Preguntas Test}}
%\input{\TemaCode/questions.tex} 



\end{document}